\chapter{Especificación de requisitos}\label{cap:especificación}

Antes de empezar con el desarrollo e implementación del sistema es necesario realizar una fase previa correspondiente
al análisis. En esta fase se van a realizar actividades que permitan extraer requisitos que va a tener el sistema con
la finalidad de tener desde el inicio un planteamiento adecuado de las funcionalidades que se tendrán. A continuación
se va a desarrollar como ha sido la fase de análisis y obtención de requisitos para el sistema

\section{Recopilación de información}

Para recopilar información se ha recurrido a la realización de entrevistas, análisis de documentos y consultas a través
de correo electrónico con el personal adecuado para el fin a realizar.

\begin{itemize}
    \item \textbf{Entrevistas}
    \newline
    Se han realizado varias entrevistas con la finalidad de ajustar el sistema lo máximo posible a lo que se necesita
    o sería interesante que hubiese en el sector laboral.

    \begin{itemize}
        \item \textbf{José (APELLIDOS) - Agricultor de (NOMBRE COOPERATIVA/EMPRESA)}
        \newline
        Con esta reunión se pudo ver cómo es el día a día de un agricultor que tiene varias fincas y cómo las gestiona.
        José mostró las distintas tablas que tiene creadas en Excel a través de las cuáles podía llevar un control exacto
        de qué actividad ha realizado en cada campaña y sobre qué finca en particular.
        Fueron varios los puntos planteados durante la reunión sobre qué podría tener la aplicación Olivaris. Entre
        ellos se destacaron:

        \begin{itemize}
            \item Funcionalidad de introducir un cuaderno de campo, siendo un punto muy interesante a tratar ya que el 
            próximo enero de 2027 va a ser obligatorio. Sobre esta funcionalidad se ha realizado un análisis más profundo
            que será comentado a continuación.
            \item Funcionalidad sobre control de costes y gastos, además de los pagos realizados a los empleados contratados.
            \item Funcionalidad sobre gestionar parcelas y las actividades realizadas en cada una de ellas, de manera que 
            se tenga una vista de todas las parcelas que son de su propiedad y que se puedan gestionar sus actividades. 
        \end{itemize}

        Gracias a este reunión se pudo obtener una aproximación de qué tipo de tablas se tendrían en la base de datos y
        como se podrían relacionar entre ellas.

        \item \textbf{Juan Luis Jiménez Laredo - Profesor titular de la ETSIIT y Director del presente TFG}
        \newline
        Las reuniones realizadas con Juan Luis han servido para llevar una revisión continua del desarrollo del proyecto,
        con la finalidad de saber qué requisitos funcionales tendrá el sistema y qué tareas hay que priorizar.
        Además, contando también con experiencia en el sector de la oliva, ha podido dar una visión del sistema más precisa,
        siendo esto de gran ayuda para definir cuáles son las funcionalidades más relevantes de la aplicación.

        Con la información extraída de la reunión realizada con Jose sobre el cuaderno de campo y con las conversaciones
        mantenidas con Juan Luis, se exploró la opción de introducir la funcionalidad del cuaderno de campo digital. 
        
        El \textbf{cuaderno de campo \cite{cuadernoCampo}} es una especie de cuaderno de bitácora que registra el uso 
        de productos químicos (fitosanitarios y fertilizantes) sobre una explotación con la finalidad de vigilar la 
        posible contaminación que pueda surgir derivada del abuso de estos productos. Este concepto está muy relacionado
        con el \textbf{Cuaderno Digital de Explotación Agrícola (CUE)}.

        Según la Conserjería de Agricultura, Pesca, Agua y Desarrollo Rural, el Cuaderno Digital de Explotación Agrícola
        (o CUE digital), es la herramienta que se va a utilizar en España para anotar las actividades que se realizan 
        en una explotación agrícola, recogiendo las actividades fitosanitarias, de fertilización o riegos aplicados.
        Para poder cumplimentar este CUE digital, es necesario que la explotación este previamente inscrita en REAFA
        (Registro Autonómico de Explotaciones Agrarias en Andalucía), ya que esta será la fuente que proporcione la 
        información de las parcelas y cultivos existentes.

        Este CUE digital, a diferencia del CUE de papel, obliga a que sea tratado desde medios electrónicos y que la
        información que lo cumplimente sea con los datos que la Administración tenga disponibles, procedentes del 
        Registro Autonómico de Explotaciones (REA).

        El aspecto que hace interesante la implementación de esta funcionalidad es que a partir del 1 de enero de 2027
        será obligatorio el uso del CUE digital para consignar la información de los tratamientos fitosanitarios. 
        
        A partir de este punto, se ha realizado una búsqueda de cómo acceder a la API que proporciona la Junta de 
        Andalucía para obtener los datos de una explotación y poder exportarlos al CUE digital. Para ello, primero es
        necesario dar de alta la "empresa" que va a desarrollar este proyecto en un portal junto con su certificado
        digital, estando aquí el primer problema: actualmente no hay constituida una empresa de desarrollo de software 
        particular, por lo que no hay certificado digital que se pueda utilizar para darse de alta y poder acceder a la
        API.

        Las soluciones buscadas han conllevado el envío de correos electrónicos a entidades que puedan suplir este problema
        y hacer de "empresa de desarrollo". Los correos enviados se van a comentar a continuación.
    \end{itemize}

    \item \textbf{Correos electrónicos}
    \newline


\end{itemize}